
$h\to\gamma\gamma$ 展示了 Higgs 如何通过带电粒子圈耦合到没有树级 Higgs 顶角的无质量规范场。把外部光子换成胶子后，同一个 Yukawa 机制给出 $h\to gg$；但这里多出非阿贝尔颜色迹与强耦合 $\alpha_s$，因而它同时连接 Higgs 扇区、夸克质量和 QCD。标准模型中只有有色费米子进入最低一圈，主导贡献来自顶夸克。

对两个在壳胶子，QCD Ward 恒等式同样把振幅限制为横向张量。沿用\eqnrefs{eq:sm-hgamma-loop-functions-dp41} 的 $A_{1/2}(\tau_q)$，可写成
\begin{equation}
\boxed{
\mathcal M^{ab}_{\mu\nu}(h\to gg)
=\delta^{ab}\frac{\alpha_s}{4\pi v_E}
\left[\sum_qA_{1/2}(\tau_q)\right]
\left(k_1\!\cdot k_2\,\eta_{\mu\nu}-k_{2\mu}k_{1\nu}\right),
\qquad
\tau_q=\frac{4m_q^2}{m_h^2}.}
\label{eq:sm-hgg-amplitude-dp94}
\end{equation}
颜色张量只有 $\delta^{ab}$，因为 Higgs 是 $SU(3)_c$ 单态，而单个夸克环的两个胶子顶角产生 $\Tr(T^aT^b)=T_F\delta^{ab}$。对标准 QCD 规范化 $T_F=1/2$，该因子已经吸收到\eqnrefs{eq:sm-hgg-amplitude-dp94} 的整体系数中。

对颜色与物理偏振求和并包含两个相同胶子的 $1/2!$，最低阶部分宽度为
\begin{equation}
\boxed{
\Gamma(h\to gg)
=\frac{G_F^{(E)}\alpha_s^2(m_hc^2)^3}
{64\sqrt2\pi^3\hbar}
\left|\sum_qA_{1/2}(\tau_q)\right|^2.}
\label{eq:sm-hgg-width-dp94}
\end{equation}
在 $m_q\gg m_h/2$ 的重夸克极限，$A_{1/2}\to4/3$，单个重夸克的贡献并不按 $1/m_q$ 消失。这种“非退耦”可直接由 Yukawa 质量关系 $m_q=v_Ey_q/(\sqrt2 c)$，即 $y_q=\sqrt2 m_qc/v_E$ 读出：环积分带来的逆质量尺度正好被 Higgs--夸克耦合中的质量因子抵消。

对顶夸克主导的低能极限，可以把重夸克积分出去。要求低能理论重现\eqnrefs{eq:sm-hgg-amplitude-dp94}，得到局域有效相互作用
\begin{equation}
\boxed{
\Lag_{hgg}^{\rm eff}
=\frac{\alpha_s}{12\pi}\frac{h}{v_E}
G^a_{\mu\nu}G^{a\mu\nu}
+\order\!\left(\frac{m_h^2}{m_t^2}\right).}
\label{eq:sm-hgg-effective-operator-dp94}
\end{equation}
这条算符也是强子对撞机中胶子融合 $gg\to h$ 的最低维描述。它说明同一个一圈匹配同时控制 Higgs 的胶子衰变和其逆过程；真正的强子截面还必须与胶子 PDF 卷积，并加入显著的高阶 QCD 修正。
